% 三、研究基础

\section{与本项目有关的研究工作积累和已取得的研究工作成绩}

\subsection{报告者本人研究工作}

\textbf{（1）紧凑坐标表示方法}

在传统表示方法下，Cosserat杆需要使用顶点的笛卡尔坐标和单元上的四元数来表示，这种表示方法使用了过多的自由度（如$7n-4$个自由度），并且需要许多额外的约束条件，包括长度约束、四元数单位化约束和轴对齐约束等。本人提出了一种紧凑坐标表示方法，摒弃了传统的节点位置自由度，转而使用基于根节点位置以及局部轴角的曲率链式参数化体系。在这种表示中，系统的求解维度被精确缩减至与其内蕴自由度完全一致的$3n$维。所有约束在表示层面自然满足，将约束优化问题转化为无约束优化问题。

\textbf{（2）近线性复杂度预处理器}

分析了系统Hessian矩阵的特殊结构，发现其具有块三对角结构。利用这一特点，设计了一个精心构造的混合预处理器。与广泛使用的块对角预处理器相比，该预处理器可以将PCG求解器的收敛速度提高一到两个数量级。更重要的是，证明了该预处理器始终是对称正定的（SPD），这为算法的稳定性奠定了坚实的理论基础。

\textbf{（3）与现有方法的对比}

与Super-Helices方法相比，本方法实现了更好的性能和稳定性，能够在相对较大的时间步长下稳定模拟具有数百个顶点和数十个碰撞的不可拉伸Cosserat杆。与RedMax方法相比，本方法专门针对细长杆优化，效率更高。

\textbf{（4）多重划分模态综合法}

观察到传统CMS方法的子结构特征模态可以视为子结构内的加窗傅里叶模态，其中边界约束固定了基函数的相位。基于这一观察，引入了从第二个空间交错划分计算得到的基函数。来自两个不同划分的子结构特征模态处于相似的频带，但它们的相位不同，从而为近似全局特征模态提供了更有效的子空间。提出了无Schur补及无局部特征求解的界面模态缩减法（SE-free IMR），利用一个划分中的界面通常位于另一个划分的子结构内部这一事实，使用子结构静力平衡问题来近似低频界面模态。实验表明，与传统CMS方法相比，本方法将精度提高了3个数量级，并在时间和内存成本上实现了更好的强扩展性和弱扩展性。

\subsection{团队研究积累}

研究团队在计算机图形学、物理仿真领域有多年积累，已在SIGGRAPH、TOG、TVCG等顶级期刊和会议发表论文多篇，积累了丰富的科研经验。承担了多项国家自然科学基金项目、国家重点研发计划项目等，具备完成高水平研究项目的能力和经验。与国外知名大学和研究机构保持密切的学术交流与合作，定期互访，及时把握国际前沿动态。

\subsection{已具备的实验条件}

硬件设施方面，实验室配备高性能计算集群（多个计算节点，支持MPI并行计算）、图形工作站（多台高性能工作站，配备专业显卡）和大容量存储（TB级存储空间，支持大规模数据管理）。软件资源方面，拥有商业仿真软件（ANSYS、Abaqus、Nastran等）、开源计算库（Eigen、PETSc、Trilinos等）和完善开发工具。文献资源方面，学校图书馆提供IEEE、ACM、ScienceDirect等数据库访问权限。

\section{已具备的实验、资料等条件，尚缺少的实验、资料条件和拟解决的途径}

\subsection{已具备条件}

硬件方面，桌面工作站配备AMD Ryzen 9 7900X 12核处理器、128GB DDR4内存、2TB SSD + 8TB HDD存储、NVIDIA RTX 3080显卡，用于算法开发、小规模测试、可视化；服务器集群配备双路Intel Xeon Gold 6248（共36核）、512GB内存，5个计算节点，InfiniBand高速互连，分布式存储系统，用于中等规模并行计算；可申请国家超算中心资源支持大规模并行计算。软件方面，操作系统采用Ubuntu 20.04 LTS / CentOS 7，开发环境包括GCC 9.3+ / Intel OneAPI、CMake 3.16+ / Make；科学计算库包括Eigen 3.4、PETSc 3.18、Trilinos 13.2、Intel MKL；图形与可视化工具包括OpenGL 4.6、VTK 9.1、ParaView 5.10；仿真与分析软件包括ANSYS Mechanical、Abaqus、MATLAB。

\subsection{尚缺少条件}

工业应用案例数据方面，缺少真实工业产品的仿真需求和数据、企业级的测试用例和验证标准、复杂工程问题的完整模型。特定领域专家指导方面，缺少航空航天领域的专家顾问、汽车工业的技术指导、高性能计算专家的支持。大规模并行计算资源方面，缺少百万核级别的超算资源、大规模GPU集群、专用的高性能存储。

\subsection{拟解决途径}

工业应用案例获取方面，建立与企业的合作关系（签署产学研合作协议、设立联合研究项目、开展技术咨询服务），参与实际工程项目（申请企业横向课题、参与企业研发项目、提供仿真技术支持），建立案例库（收集公开的工业案例、整理测试数据集、建立验证标准）。专家支持获取方面，聘请校外专家作为指导小组成员、定期组织专家咨询会议、邀请专家做学术报告；参加国内外学术会议、参与学术研讨班、组织专题研讨会；加入相关学术组织、参与学术联盟、建立长期合作关系。计算资源获取方面，申请国家超级计算中心资源、参与国家超算项目、利用高校超算联盟资源；利用阿里云、腾讯云等平台、申请教育优惠、按需使用弹性资源；与其他高校共享计算资源、参与计算资源共享计划、建立校际合作机制。

\subsection{资源保障计划}

经费保障方面，有国家自然科学基金项目经费、学校研究生创新项目资助、实验室运行经费支持。人员保障方面，研究团队稳定，研究生梯队建设完善，技术支持人员配备齐全。时间保障方面，四年博士学制，全日制研究，合理的研究计划安排。